
\begin{figure}[h]
	\centering
        \includegraphics[width=0.84\linewidth]{figs/miniF2F_SOTA.pdf}
		\caption{Growth in MiniF2F-Test performance over time.}
		\label{fig:miniF2F_SOTA}
	\vspace{-0.5in}
\end{figure}

\section{Introduction}

Recent advances in large language models (LLMs) have shown extending reasoning lengths through natural language can significantly boost performance on math benchmarks such as MATH and AIME \citep{openai2024openaio1card,deepseekai2025deepseekr1incentivizingreasoningcapability}. 
Training such models requires reinforcement learning (RL) on verifiable answers. 
It is extremely difficult to automatically, or even manually, verify a proof in natural language since each step must be carefully checked for correctness \citep{petrov2025proofbluffevaluatingllms}.
This poses significant challenges for applying reinforcement learning to the training of large language models to prove mathematical statements.
Unlike natural language, formal languages such as Lean can provide a clear and automatic signal on the correctness of a formalized proof.    % which is suitable for RL.
A noteworthy work from AlphaProof \citep{alphaproof} uses Lean to successfully solve 3 problems from the 2024 International Mathematical Olympaid (IMO).
AlphaProof demonstrates that LLMs using formal language are capable of proving very challenging problems that LLMs using natural language fail to prove.

There are two types of LLM formal provers, step-level provers \citep{alphaproof,polu2022formal,wu2024LeangithubcompilinggithubLean,wu2024internlm25stepproveradvancingautomatedtheorem,xin2025bfsproverscalablebestfirsttree} and whole-proof generation provers\citep{xin2025deepseekproverv,dong2025stpselfplayllmtheorem}.
Step-level provers incrementally generate Lean code line-by-line. While this enables close interaction with the Lean environment, it requires special scaffolding to generate a complete Lean proof, and the interaction is often too granular to allow high-level reasoning.
In contrast, whole-proof models generate an entire Lean proof at once, but typically lack interaction with the Lean compiler.
Recent work has shown that combining whole-proof models with long chain-of-thought reasoning \citep{kimina,deepseek-v2,lin2025goedel} substantially outperforms step-level provers.
% interact more
In this work, we propose \textbf{Seed-Prover}, a whole-proof model with following features:
\begin{itemize}
    \item \textbf{Lemma-Style Proving}: Seed-Prover tries to generate useful intermediate lemmas before proving the main statement. These lemmas serve as shared knowledge across different inference paths.
    \item \textbf{Iterative Proof Refinement}: Seed-Prover iteratively refines its proof based on Lean compiler feedback, previous proved lemmas, and self-summarization.
    \item \textbf{Test-Time Scaling}: We implement a three-tiered inference strategy that enables Seed-Prover to think both deeply and broadly—allocating thinking budget to fine details while exploring interesting properties.
    \item \textbf{SOTA Performance}: Seed-Prover proves 5 out of 6 problems in IMO 2025, saturates MiniF2F \citep{zhengminif2f} (shown in Figure~\ref{fig:miniF2F_SOTA}), and outperforms prior work by up to $3\times$ on multiple formal benchmarks.
\end{itemize}

% Seed-Geometry
Due to the lack of sufficient geometry support in Lean, Seed-Prover incorporates a dedicated geometry reasoning engine \textbf{Seed-Geometry}. Similar to existing line of efforts in AlphaGeometry \cite{trinh2024solving,chervonyi2025gold} and TongGeometry \cite{zhang2024proposing}, Seed-Geometry follows the forward-chaining design in the reasoning engine implementation, where the system derives all known facts by checking applicable rules until closure is reached. By backward-tracing fact dependencies, Seed-Geometry identifies the minimum dependency relations in a geometry problem's configuration, seperating the problem context from the auxiliary constructions necessary to prove a problem. Using statistics derived from more than past 20 years of math olympiad competitions, Seed-Geometry performs extensive search in the geometry space defined by its dedicated domain-specific language and establishes a repository of 230 million unique geometry problems requiring auxiliaries. A Seed model trained on such dedicated geometry data becomes an exceptionally effective neuro-symbolic geometry prover, where it fills in the missing auxiliary geometry elements and the geometry reasoning engine performs step-by-step forward-chaining, completing the final proof of a problem. In experiments, Seed-Geometry solves 43 of the IMO-AG-50 (\textit{vs.} 42 by AlphaGeometry 2), a benchmark that curates geometry problems of IMO from 2000 to 2024. It also sets a new state-of-the-art on the IMO shortlist geometry problems from 2000 to 2022, and notably solves the geometry problem of IMO 2025 under just 2 seconds.